\chapter{What is Security Needed for?}

\section{Introduction}
The security services that will enable widespread use and acceptance of
commerce over the Internet include the following:

\begin{htmllist}
  \htmlitemmark{RedBall}
  \item[Authentication:] whereby an individual, an organization, or a
    computer can prove its identity.
  \item[Authorization:] the ability of a system, once identity is verified,
    to control access to specific resources.
  \item[Confidentiality:] the ability to maintain the secrecy of the
    contents of a transmission between authorized parties.
  \item[Integrity:] the capability of ensuring that a transmission arrives
    at its destination in exactly the same form as it was sent.
  \item[Non-repudiation of origin:] the ability to ensure that when an
    entity sends an authenticated electronic communication, it cannot later
    deny the origin, or contents of the communication.
\end{htmllist}

In applications of electronic commerce, encryption is typically used to
provide security. However, the industry seems to recognize that
cryptography, while a critical component of secure networking, is not
enough to protect a business from fraud. Cryptography provides for
transaction security, but does not do much to prevent unauthorized
access to information and accounts. A good analogy would be using an
armored van to transport cash from one branch to another, and then
leaving it in the middle of the lobby once it reaches its destination.
Information needs to be ``put in the vault'' once it reaches its
destination, and another component is needed to do this: a firewall to
secure this local network from Internet access.\\

Security methods, such as Internet firewalls, are very popular now, but
many organizations may believe, or be led to believe, that an Internet
firewall alone is sufficient for securing their network. It is like
getting the most secure front door money can buy for your house but
leaving the garage door unlocked, or the same, weak sliding door
entrance from your back deck. Only if everyone plays by the same rules
is it effective. Thus, we may expect to see smart cards and their
related technology flourishing, as a mean to provide businesses with
reliable access control and authentication processes: smart cards may
be used to carry and to prove the user's identity, as well as other
relevant personal information, and all this in a secure manner.\\

However, firewalls and smart cards are not enough. A risk and business
analysis is almost always required, leading into the development of a
security policy and the prescription of security mechanisms and methods
for implementation within the local computing environment. Besides,
doing this once is not enough. Threats change, vulnerabilities change,
business requirements change, and the available counter-measures
change: all of these must be periodically and routinely reevaluated.\\


\section{Local Computing Environment}

\subsection{Introduction}
Adequately secure systems deter, prevent, or detect unauthorized
disclosure, modification, or use of information. Company information
requires protection from intruders, as well as from employees with
authorized computer access privileges who attempt to perform
unauthorized actions. Protection is achieved not only by technical,
physical and personnel safeguards, but also by clearly articulating and
implementing company policy regarding authorized system use to
information users and processing personnel at all levels.

\subsection{The risks}
The spread of computing power to individual employees via personal
computers, local-area networks (LAN), and distributed processing has
drastically changed the way we manage and control information
resources. Internal controls and control points that were present in
the past when we were dealing with manual or batch processes have not
been established in many of today's automated systems. Reliance upon
inadequately controlled computer systems can have serious consequences,
including:

\begin{itemize}
  \item Inability or impairment of the company's ability to perform its
    mission
  \item Inability to provide needed services to the public
  \item Waste, loss, misuse, or misappropriation of funds
  \item Loss of credibility or embarrassment to the others
\end{itemize}

To avoid these consequences, a broad set of information security issues
must be effectively and comprehensively addressed.

\subsection{The goals}

\begin{htmllist}
  \htmlitemmark{RedBall}

    \latexhtml{\item {\bf Achieve cost-effective security}}{\item[Achieve
    cost-effective security]}
    
    The dollars spent for security measures to control or contain losses
    should never be more than the projected dollar loss if something
    adverse happened to the information resource. Cost-effective security
    results when reduction in risk through implementation of safeguards is
    balanced with costs. The greater the value of information processed, or
    the more severe the consequences if something happens to it, the
    greater the need for control measures to protect it. The person who can
    best determine the value or importance of data is the functional
    manager who is responsible for the data. For example, the manager
    responsible for the company's budget program is the one who should
    establish requirements for the protection of the automated data which
    supports the program. This manager knows better than anyone else in the
    organization what the impact will be if the data is inaccurate or
    unavailable. Additionally, this manager usually is the supervisor of
    most of the users of the data.\\
    
    It is important that these trade-offs of cost versus risk reduction be
    explicitly considered, and that management understand the degree of
    risk remaining after selected controls are implemented.

    \latexhtml{\item {\bf Assure operational continuity}}{\item[Assure
      operational continuity]}
    
    With ever-increasing demands for timely information and greater volumes
    of information being processed, the threat of information system
    disruption is a very serious one. In some cases, interruptions of only
    a few hours are unacceptable. The impact due to inability to process
    data should be assessed, and actions should be taken to assure
    availability of those systems considered essential to agency operation.
    Functional management must identify critical computer applications and
    develop contingency plans so that the probability of loss of data
    processing and telecommunications support is minimized.

    \latexhtml{\item {\bf Maintain integrity}}{\item[Maintain integrity]}
    
    Integrity of information means you can trust the data and the processes
    that manipulate it. Not only does this mean that errors and omissions
    are minimized, but also that the information system is protected from
    deliberate actions to wrongfully change the data. Information can be
    said to have integrity when it corresponds to the expectations and
    assumptions of the users.

    \latexhtml{\item {\bf Assure confidentiality}}{\item[Assure confidentiality]}
    
    Confidentiality of sensitive data is often, but not always, a
    requirement of company systems. Privacy requirements for personal
    information is dictated by statute, while confidentiality of other
    company information is determined by the nature of that information,
    e.g., information submitted by the company to its resellers. The impact
    of wrongful disclosure must be considered in understanding
    confidentiality requirements.
\end{htmllist}

\subsection{Control decisions}
The official responsible for the company function served by the
automated information system has a critical role in making decisions
regarding security and control. In the past, risk was often
unconsciously accepted when such individuals assumed the computer
facility operators were taking care of security. In fact, there are
decisions to be made and security elements to be provided that cannot
be delegated to the operator of the system. In many cases, the user or
manager develops the application and operates solely.\\

The cost of control must be balanced with system efficiency and
usability issues. Risk must be evaluated and cost-effective controls
selected to provide a prudent level of control while maximizing
productivity. Controls are often closely connected with the system
function, and cannot be effectively designed without significant
understanding of the process being automated.

\section{What Does the Use of the Web Imply About Security?}
As is often the case with systems that provide a high degree of data
sharing, there are a number of security issues for sites that want to
take advantage of the web. They can be broken down into three major
categories that will be exposed below.

\subsection{Client threats}
The major security issue with running client software comes about given
the nature of the web: client programs interpret data that is
downloaded from arbitrary servers on the Internet. If there are no
checks on the contents of this imported data, there exists the
potential for this data to subvert programs running on the client
systems. These ``Trojan horses'' may take several forms, from malicious
URLs to rogue code that is run through interpreters (such as
PostScript) on the client system.\\

Another significant security issue applies mostly to sites that are
protected by Internet firewalls. These sites have implemented security
policies in their firewalls that describe the network services they
want to permit within their organizations. A feature of URLs is that
they support several resource types: file, Ftp, Http, Gopher, Wais,
Nntp, Mailto, Prospero, Telnet, and Rlogin (other resources are in the
works). It is likely that the firewall implementation permits only a
subset of these resources. Therefore, since web servers provide these
services independently of their normal mechanisms (independent of the
normal Ftp channels, for example), it is possible to bypass a firewall
security policy by using web services.

\subsection{Server threats}
Threats to servers consist broadly of unauthorized modification of
server data, authentication and access control to server data, and
compromising a server system by exploiting bugs in the web server
software.\\

The HTTP protocol provides several methods (for example, PUT) for
writing data on a server. It is possible that some of these methods may
be used for unauthorized modification of server data. Work is ongoing
to address security concerns in HTTP servers.

\subsection{Point to point security issues}
There are a number of security issues that arise between clients and
servers. These security requirements depend on each other in many cases:

\begin{htmllist}
  \item[Authentication:] For some potential uses of the web (e.g.,
    electronic commerce), it is important that clients authenticate
    themselves to servers, that servers authenticate to clients, and that
    both authenticate to each other. It is important that a form of
    authentication be used that is not easily spoofed, such as
    cryptographically signed certificates.
  \item[Confidentiality:] Confidentiality is important for web uses that
    involve sensitive data. Data exported by web clients or servers may
    need to be protected from eavesdropping. For example, services are
    available (on-line florist) that require the passing of credit card
    information. This requirement will be amplified when other kinds of
    data, such as employee records, government files, etc., begin
    traversing the web.
  \item[Integrity:] Certain transactions have the requirement that
    transaction requests and/or contents remain unmodified during delivery.
    The electronic commerce example is a case where data integrity is
    critical to proper functioning of the system.
  \item[Availibility:] This means all the systems are up and running when
    they are needed.
\end{htmllist}
